\chapter{Agent reliability and verification engineering}

\section{Engineering principles}

Reliability is a property of the entire workflow. Models may hallucinate completion, repeat failing actions, exhaust context or use a tool incorrectly. External state records prevent a restart from silently repeating side effects. Timeouts and resource budgets protect infrastructure; idempotent operations and checkpoints reduce partial-failure risk.\src{A9,O7}

\section{Implementation approach}

Classify failures as transient infrastructure, deterministic test failure, policy denial or ambiguous specification. Retry transient failures with bounded backoff; repair test failures only when scope permits; escalate denials and ambiguity. Persist task state before and after consequential transitions. On resume validate baseline and candidate identity. Treat missing verifier output as a failure and never infer success from an exit code alone when evidence is required.

\section{Verification and limitations}

Inject failures such as a timeout, nonzero exit, stale commit and missing report. Assert that each ends in an auditable terminal state. The sample orchestrator demonstrates bounded attempts, not a durable distributed scheduler.

\section{Chapter conclusion}

Apply the stated verification criteria before expanding the scope of this workflow. Record deviations, evidence and unresolved risks in the task record.
